\section{Preliminary Results}
\label{sec:eval}

\paragraph{Simulation setup.}
\lijl{For the workshop paper, this part could be condensed.}
We evaluate \dht using a simulated DSN.
Unless otherwise specified, the simulations involve 12,000 nodes, each with a storage capacity ranging from $2^{10}$ data units (minimum) to $2^{17}$ data units (maximum), resulting in node class values between 10 and 17.
The distribution of node capacities follows a shifted Zipf distribution~\cite{zipf} with a skewness parameter of 1.
In addition to \dht, we simulate a standard Kademlia DHT as a baseline for comparison.
Each piece of content is replicated three times, i.e., $r=3$.
Finally, each simulation is run 100 times, and we report the average and standard error of the results.

\begin{figure}
    \centering
    \includegraphics[width=0.7\linewidth]{data/simulation/graphs/util}
    \caption{Effective overall utilization achieved by DSNs.}
    \label{fig:util-loss}
    \Description{}
\end{figure}

Initially, we conducted 10 million storage attempts with uniformly random keys on the simulated DSNs.
With \dht, all storage attempts were successful, and no node exhausted its storage space.
The outcomes for Kademlia are illustrated in \autoref{fig:overwhelm} and detailed in \autoref{sec:bg:problem}.

Next, we conducted a stress simulation by continuously storing data with random keys (with random eviction upon node storage exhaustion) until data loss occurred, and we report the overall utilization in \autoref{fig:util-loss}.
We performed simulations with varying skewness of node capacities, ranging from 1 (highly diverse capacities) to 2 (where almost all nodes have the same $2^{10}$ capacity).
We also conducted an ablation study to assess the effectiveness of the classified distance and 2RC designs independently.
The results indicate that both designs contribute to improved utilization, but in different scenarios.
Classified distance is particularly effective when node capacities are diverse (low skewness), but its performance approaches that of Kademlia as the skewness increases and most nodes converge to the same (lowest) node class.
Conversely, 2RC achieves near-perfect utilization under high skewness but becomes less effective when capacity heterogeneity is significant.
Ultimately, \dht combines the benefits of both classified distance and 2RC, outperforming Kademlia by over 60\% in utilization across all skewness levels and reaching almost 100\% utilization at high skewness.
These simulation results demonstrate that \dht effectively achieves its goal of optimizing effective overall utilization, and both its constituent designs perform as expected.

We proceed with micro-benchmarks to provide an in-depth analysis of \dht, detailed in appendix due to the lack of space.
To summarize, the classified distance can effectively balance the probability of utilizing storage units on every node, and the captured available space snapshots prove that \dht perform the equalization well.